Ce projet étudie le choix des stations et le tracé d'une ligne de transport circulaire
au moyen du problème Ring-Star, ou anneau-étoile. Parmi $n$ points candidats,
exactement $p$ deviennent des stations, dont un sommet imposé. Les stations sont
reliées par un cycle unique et chaque autre point est affecté à une station.
Le coût additionne la longueur du cycle, pondérée par $\alpha$, et les distances
d'accès aux stations. Le modèle représente une décision de localisation et de
routage ; il ne constitue pas une étude opérationnelle d'un réseau de transport réel.

Le travail reprend l'énoncé de la SAÉ Optimisation de Sup Galilée et en explicite
les hypothèses. Lorsque $p=n$ et $\alpha>0$, le problème contient le voyageur de
commerce euclidien comme cas particulier, ce qui établit sa NP-difficulté.
La résolution exacte repose sur une programmation linéaire en nombres entiers
avec des variables d'affectation, des arêtes d'anneau et un flot assurant la
connexité. Le modèle est construit avec PuLP et résolu par CBC.

Trois méthodes approchées complètent cette formulation. Une construction
déterministe choisit les stations à partir d'une grille, affecte les points à leur
station la plus proche et construit un cycle par plus proche voisin, amélioré
par 2-opt. Une recherche locale accepte les échanges de stations améliorants.
Une recherche tabou peut accepter une dégradation temporaire, interdit certains
retours et conserve la meilleure solution rencontrée.

L'évaluation principale porte sur huit nuages de points issus de TSPLIB, de 16 à
100 sommets, avec $p\in\{3,5,10\}$ et $\alpha=1$. Les deux recherches sont répétées
sur trois graines. Le protocole compte 192 cas de méthode et de graine :
168 exécutions heuristiques, quatre résolutions PLNE et 20 cas PLNE ignorés
par les limites de taille configurées. Sur les quatre PLNE exécutés, trois
aboutissent à une preuve d'optimalité et un fournit une solution réalisable
sans preuve dans le budget de cinq secondes.

Les gains médians de coût face à la construction sont de 11,79\,\% pour la
recherche locale et de 13,67\,\% pour la recherche tabou, en donnant le même poids
à chacun des 24 cas d'instance et de cardinalité. La moyenne tabou est inférieure
à la moyenne locale dans 22 cas et égale dans deux cas. Ce constat est lié au
protocole : une itération tabou évalue jusqu'à 15 échanges, contre un seul pour
la recherche locale. Il ne démontre donc pas une supériorité à effort de calcul égal.

L'étude détaillée d'\code{att48}, avec dix stations, illustre une baisse du coût
de 48\,833,89 à 42\,743,75, soit 12,47\,\%. La reproductibilité repose sur la
conservation des coordonnées, des paramètres, des solutions, des traces et des
versions logicielles. Les 172 solutions de la campagne ont été relues et leurs
coûts recalculés pour le rapport. La suite de 101 tests passe également avec
un contrôle renforcé de l'export des cas ignorés.

\vspace{3mm}
{\sffamily\small\raggedright\textbf{Mots-clés :} recherche opérationnelle, Ring-Star,
localisation-routage, PLNE, recherche tabou, expérimentation reproductible.\par}
